\documentclass[10pt, twocolumn, a4paper]{article}

\usepackage[margin=0.65in]{geometry}
\usepackage{amsmath, amssymb, amsfonts}
\usepackage{booktabs}
\usepackage{graphicx}
\usepackage{hyperref}
\usepackage{microtype}
\usepackage{abstract}
\usepackage{cite}
\usepackage{array}

\hypersetup{
    colorlinks=true,
    linkcolor=blue,
    citecolor=blue,
    urlcolor=blue
}

\renewcommand{\abstractnamefont}{\normalfont\bfseries}
\renewcommand{\abstracttextfont}{\normalfont\small}

\title{\LARGE \bfseries Linear Steering Fails to Calibrate Peer Influence: Evaluating A-LQR Control in a Two-Agent Resist/Accept Task}

\author{
    Manu Pareek\\
    Georgia Institute of Technology\\
    \texttt{mpareek3@gatech.edu}
}

\date{\today}

\begin{document}

\makeatletter
\twocolumn[
  \begin{@twocolumnfalse}
    \maketitle
    \begin{abstract}
    Closed-loop representation control methods such as Activation Linear Quadratic Regulator (A-LQR) steering have shown success modulating single-agent behavioral traits (e.g., toxicity, truthfulness, refusal) \cite{skifstad2026alqr}. Whether such control can calibrate \emph{social} susceptibility---resisting bad peer pressure while accepting good corrections---is untested. We built a two-agent resist/accept harness on TruthfulQA MC1 with \texttt{gemma-2-2b-it} and extracted contrastive A-LQR steering vectors for the resist and accept objectives \cite{pareek2026multiagent}. Unsteered, the agent exhibits acute overinfluence: it answers most resist items correctly alone yet retains the correct answer after an assertive wrong peer in only 1/10 cases. Critically, A-LQR steering never improves resistance across $\lambda \in {[0.5, 12.0]}$---mild gains leave resistance unchanged, moderate gains eliminate it, and high gains collapse decoding entirely. On a pilot sweep, steering on several items moved the agent \emph{toward} the incorrect peer. We conclude that peer compliance is not a linearly steerable direction in this regime: the coupling between instruction-following and social deference resists single-direction control, and calibrated influence likely requires non-linear or gated interventions.
    \end{abstract}
    \vspace{0.8em}
  \end{@twocolumnfalse}
]
\makeatother

\section{Introduction \& Motivation}
Representation-engineering methods allow direct intervention on a Large Language Model's (LLM) internal activations at generation time without modifying weights \cite{skifstad2026alqr}. A-LQR steering in particular models hidden-state dynamics as a discrete-time linear system and computes closed-loop feedback toward target representations derived from contrastive prompt pairs \cite{skifstad2026alqr}.

Single-agent steering has proven effective for scalar behavioral traits---sentiment, toxicity, refusal. Multi-agent interaction introduces a harder target: \emph{calibrated} social influence, where an agent must condition its compliance on the correctness of the peer rather than the assertiveness of the message \cite{pareek2026multiagent}:
\begin{itemize}
    \item \textbf{Overinfluence / Sycophancy}: Agent B abandons correct knowledge to agree with a confident, incorrect peer (Agent A).
    \item \textbf{Underinfluence / Resistance Failure}: Agent B rejects valid corrections from Agent A and persists in error.
\end{itemize}

The central question: \textit{can contrastive activation steering decouple an agent's compliance behavior from the correctness of what it is complying with?} If peer deference were a separable linear direction, A-LQR should reduce overinfluence while preserving receptivity to good corrections. Our results indicate it is not.

\section{Experimental Setup \& Methodology}
We established a two-agent simulation environment \cite{pareek2026multiagent}:

\subsection{Agent Architecture \& Stack}
\begin{itemize}
    \item \textbf{Target Agent (Agent B)}: \texttt{google/gemma-2-2b-it} via HuggingFace, instruction-tuned, official chat templates \cite{gemma2024}.
    \item \textbf{Peer Agent (Agent A)}: A scripted peer delivering assertive feedback to Agent B.
    \item \textbf{Task}: TruthfulQA MC1, exact-letter scoring \cite{lin2021truthfulqa}. Constraining evaluation to exact-letter selection removes judge noise from open-ended generation.
\end{itemize}

\subsection{Behavioral Objectives}
\begin{enumerate}
    \item \textbf{Resist}: Agent B answers correctly alone. Agent A asserts a wrong choice (e.g., \textit{``I'm sure it's C. Reply with only: C''}). Success = B keeps its correct answer.
    \item \textbf{Accept}: Agent A asserts the correct choice. Success = B outputs it.
\end{enumerate}
Both conditions are scored identically (correct letter after the peer message), isolating \emph{who suggested} the answer rather than \emph{what} the answer is.

\subsection{Contrastive Sets \& A-LQR Vector Extraction}
Steering directions and Jacobian matrices were computed across target layers to form control setpoints:
\begin{equation}
\Delta X_{\text{resist}} = \mathrm{acts}(\mathrm{resist}) - \mathrm{acts}(\mathrm{defer})
\end{equation}
\begin{equation}
\Delta X_{\text{accept}} = \mathrm{acts}(\mathrm{accept}) - \mathrm{acts}(\mathrm{ignore})
\end{equation}
Contrastive sets held the question and peer message fixed while varying Agent B's instruction (e.g., \textit{``only use A if correct''} vs.\ \textit{``always agree with A''}). Pickles were rebuilt with matched chat templates after earlier plain-text extractions introduced noise. A-LQR control was applied to B strictly \textbf{after} processing A's message.

\subsection{Influence Measurement}
Beyond binary resist/accept counts, the harness records a continuous influence score per item: the shift in B's output toward A's position after exposure to the peer, i.e.\ $\mathrm{sim}(B_{\mathrm{after}}, A) - \mathrm{sim}(B_{\mathrm{alone}}, A)$ \cite{pareek2026multiagent}.

\section{Empirical Results}

\subsection{Baseline: Overinfluence Is Real, Not Ignorance}
The strongest evidence that this is a social-compliance failure rather than a capability gap is the alone-vs-peer gap on resist items (Table~\ref{tab:main}): Agent B produced the correct letter on the majority of resist items when answering alone, yet after the assertive wrong peer retained the correct answer in only \textbf{1/10} cases. On accept items, the peer was followed in \textbf{9/10} cases---baseline compliance is near-ceiling in both directions.

\subsection{Steering Sweeps}
Table~\ref{tab:main} summarizes the main sweep on TruthfulQA MC1 (10 resist + 10 accept items per condition):

\begin{table*}[t]
\centering
\caption{Resist/Accept performance on TruthfulQA MC1 across conditions ($n=10$ per objective)}
\label{tab:main}
\vspace{0.5em}
\begin{tabular}{lcccp{5.6cm}}
\toprule
\textbf{Condition} & \textbf{All} & \textbf{Resist} & \textbf{Accept} & \textbf{Generation quality} \\
\midrule
Alone (no peer, unsteered) & --- & $\sim$6/10 & --- & Pre-peer capability is mostly intact \\
Peer, unsteered (baseline) & 10/20 & 1/10 & 9/10 & Fluent; acute overinfluence \\
A-LQR $\lambda \in {[0.5, 1.5]}$ & 11/20 & 1/10 & 10/10 & Fluent; no resistance gain \\
A-LQR $\lambda \in {[2.0, 5.0]}$ & 10/20 & 0/10 & 10/10 & Letter spam (``C C C\ldots''); zero resistance \\
A-LQR $\lambda \in {[8.0, 12.0]}$ & 0/20 & 0/10 & 0/10 & Severe decoding collapse \\
\bottomrule
\end{tabular}
\end{table*}

\subsection{Ablation Across Pilot Datasets}
The pattern replicated across earlier pilots before the matched TQA-MC rebuild \cite{pareek2026multiagent}:
\begin{itemize}
    \item \textbf{Handmade pilot ($n=20$, $\lambda=1.25$)}: overall accuracy $14/20 \to 11/20$; resist $6/11 \to 3/11$; accept unchanged $8/9 \to 8/9$. Notably, on several resist items steering \emph{increased} movement toward the incorrect peer---the push landed on the wrong side of the compliance direction.
    \item \textbf{No-Country sweep ($n=18$)}: resist declined monotonically with gain, from $5/10$ at baseline to $0/10$ at $\lambda=2.0$; accept held near baseline until degradation set in.
\end{itemize}

\section{Discussion \& Key Findings}
\begin{enumerate}
    \item \textbf{Compliance is not a linearly steerable direction.} Contrastive vectors cleanly shift single-agent traits, but peer-induced deference appears entangled with instruction-following and the token distribution itself. Steering toward ``resist'' did not selectively reduce overinfluence---on several items it increased deference to the wrong peer.
    \item \textbf{Degradation horizon precedes any benefit.} Increasing $\lambda$ never produced a regime of improved resistance. Because A-LQR scales setpoint strength as $\beta = \lambda \times \|E\|$ per layer, large gains warp hidden states off the natural generation manifold (letter spam at $\lambda \geq 2$; full collapse at $\lambda \geq 8$) before influencing social behavior.
    \item \textbf{Asymmetric headroom.} Accept is already near-ceiling at baseline (9/10), so nearly all measurable harm concentrates on resist---where steering was consistently neutral-to-negative. The $\lambda$ knob is ruled out across 0.5--12; this is a method limitation, not a tuning gap.
    \item \textbf{Template matching matters.} Mismatches between plain-text pickles and chat-template evaluation introduced noise; matched templates stabilized extraction but did not change the qualitative outcome.
\end{enumerate}

\section{Limitations \& Future Directions}
\begin{itemize}
    \item \textbf{Scale}: $n=20$ per condition on one 2B-parameter model with a single seed limits statistical strength; the consistent replication across three dataset/model variants mitigates but does not eliminate this.
    \item \textbf{Scripted peer}: A was a static assertive template. Softer or multi-turn peers, and a live (unscripted) Agent A, remain unevaluated---deliberately, until steered resist beats baseline.
    \item \textbf{Missing prompt baseline}: A prompt-engineering control (e.g., ``ignore the peer if unsure'') was identified but not yet run; it would establish whether the failure is specific to activation-level methods.
    \item \textbf{Future methods}: Gated or nonlinear control (steering conditioned on detected sycophancy), richer influence pickles, and multi-agent SFT are natural successors on the same harness.
\end{itemize}

\section*{Acknowledgments}
This work was carried out as part of a research project with the Trustworthy Robotics Lab at the Georgia Institute of Technology and builds directly on the lab's A-LQR activation-steering framework \cite{skifstad2026alqr}. The author thanks the lab's graduate researchers for mentorship and infrastructure support.

\begin{thebibliography}{99}

\bibitem{skifstad2026alqr}
A. Skifstad et al., ``Closed-Loop Control of Language Model Activation States via Linear Quadratic Regulation (A-LQR),'' \textit{arXiv preprint arXiv:2604.19018}, 2026.

\bibitem{lin2021truthfulqa}
S. Lin, J. Hilton, and O. Evans, ``TruthfulQA: Measuring How Models Mimic Human Falsehoods,'' \textit{arXiv preprint arXiv:2109.07958}, 2021.

\bibitem{gemma2024}
Gemma Team, ``Gemma 2: Improving Open Language Models via Two-Stage Knowledge Distillation,'' \textit{Google DeepMind Technical Report}, 2024.

\bibitem{pareek2026multiagent}
M. Pareek, ``Multi-Agent Calibrated Influence with A-LQR: Resist/Accept Peer Pressure on TruthfulQA-MC,'' \textit{GitHub Repository}, 2026. \url{https://github.com/manupareekk/multiagent-alqr-influence}

\end{thebibliography}

\end{document}
